This thesis presents methods for the valuation and discretization of barrier options, financial derivatives whose payoff depends on whether the underlying asset price crosses a predetermined barrier level during the option's lifetime. Barrier options are widely used in financial markets for risk management and investment strategies due to their lower cost compared to standard options.

The main objectives of this work are to analyze and implement methods for pricing barrier options, with particular emphasis on discretization techniques that accurately capture the barrier condition. The methodology includes four main approaches: (1) analytical formulas under the Black-Scholes model, which assume a lognormal behavior of the asset price and provide exact closed-form solutions for European options; (2) Monte Carlo simulation, a versatile and robust computational method to approximate the value by generating multiple random price trajectories; (3) advanced discrete models, specifically the Adaptive Mesh Model (AMM), which addresses the positioning error problem when the barrier does not align with standard tree nodes; and (4) static replication, a technique for hedging barrier options by constructing a portfolio of vanilla options that reproduces the same payoff under certain conditions. We focus on both European and American-style barrier options.

The results demonstrate the effectiveness of the proposed discretization schemes in achieving accurate valuations while maintaining computational efficiency. We compare the performance of different numerical approaches and analyze their convergence properties. The work contributes to the understanding of numerical methods for exotic options pricing and provides practical tools for financial risk assessment.

\textbf{Keywords:} barrier options, option pricing, Black-Scholes model, Monte Carlo simulation, binomial and trinomial trees, static replication, adaptive meshes, financial derivatives.
